%!TEX root = forallx.tex
\thispagestyle{empty}
{\Huge\forallx}

{\tt Introdução à Lógica Formal}

\vfill

{\sf P.D. Magnus}\\
\emph{University at Albany, State University of New York}

\vfill

{\sf \textbf{Tradução para o português brasileiro}}\\
{\sf Carlos André Duarte Costa}\\
{\sf 26 de setembro de 2026}

\vfill

{\sf
  \href{https://www.fecundity.com/logic/}{fecundity.com/logic}, versão 1.4 [\bookversion]\\
  Este livro é oferecido sob uma licença Creative Commons.\\
  (Atribuição 4.0)
}

\newpage
\thispagestyle{empty}%

{\sf
O autor gostaria de agradecer às pessoas que tornaram este projeto possível. Entre elas, destacam-se Cristyn Magnus, que leu muitas versões preliminares; Aaron Schiller, um dos primeiros a adotar o livro, que forneceu comentários extensos e úteis; e Bin Kang, Craig Erb, Nathan Carter, Wes McMichael, Selva Samuel, Dave Krueger, Brandon Lee, Toan Tran, Marcus Adams, Matthew Brown e os estudantes da disciplina de Introdução à Lógica, que identificaram vários erros em versões anteriores do livro.
}

\vfill
{
\copyright\ \ifthenelse{\year=2005}{\number\year}{2005--\number\year} por P.D. Magnus. Alguns direitos reservados.\\
\copyright\ 2026 Carlos André Duarte Costa (tradução para o português brasileiro).
}

{\footnotesize
Você tem liberdade para copiar este livro, distribuí-lo, exibi-lo e criar obras derivadas, sob a seguinte condição: Atribuição. Você deve creditar o autor original. --- Em qualquer reutilização ou distribuição, deve deixar claros a outras pessoas os termos da licença desta obra. Qualquer uma dessas condições pode ser dispensada mediante autorização do titular dos direitos autorais. Seus direitos de uso justo e demais direitos não são afetados pelo exposto acima. --- Este é um resumo legível por humanos da licença completa, disponível on-line em \url{http://creativecommons.org/licenses/by/4.0/}
}

{
A diagramação foi realizada inteiramente em \LaTeX$2\varepsilon$. O estilo de diagramação das provas baseia-se em fitch.sty (v0.4), de Peter Selinger, da University of Ottawa. Esta tradução foi preparada a partir da versão \editionoriginalversion\ da obra original. A versão mais recente da obra original está disponível on-line em \url{https://www.fecundity.com/logic/}.
}

\newpage
\thispagestyle{empty}%

{\Large\sf\textbf{Créditos do projeto REALMat}}

\bigskip

{\sf
Projeto REALMat — Recursos Educacionais Abertos na Licenciatura em Matemática.\par
\medskip
O REALMat organiza e disponibiliza esta edição brasileira.\par
\medskip
\textbf{Tradução para o português brasileiro}\par
\smallskip
Carlos André Duarte Costa\par
\medskip
\textbf{Revisão editorial da edição brasileira}\par
\smallskip
Equipe de revisão da edição brasileira\par
\bigskip
Portal: \url{https://projetorealmat.github.io/}\par
\smallskip
Repositório desta edição: \url{https://github.com/projetorealmat/forallx}\par
\smallskip
Fonte original: \url{https://github.com/OpenLogicProject/forallx}.
}